\begin{papercard}{doan2025}
\cf{Stages} \stagetag{A} (purposive selection of six companies), \stagetag{D} (dated policy capture), \stagetag{C} (deductive CI coding plus a bottom-up codebook), \stagetag{I} (single coder with consultation, flagged as a reliability gap), \stagetag{H} (stated risks and benefits, collective harms). The study analyses the EU/UK privacy and research-consent policies of six direct-to-consumer (DTC) genetic-testing companies with Nissenbaum's contextual integrity (CI) flow parameters, vague-term lists, GDPR transparency articles and a user-relevance attribute set. Its coding design is the closest template in the corpus for a policy-text analysis by one researcher.

\cf{Research questions} RQ1: how transparent are the described information flows in terms of including all parameters, using vague language, and parameter bloating (more than two concepts in one parameter)? RQ2: how aligned are the flows with GDPR transparency requirements? RQ3: how aligned with user-relevant information categories for research data sharing (data user, data collector, reason for data use, information and consent, ethical review)? RQ4: how many and which risks and benefits are stated? \pc{doan2025}{5}{RQ1: How transparent are information flows described in the available policies in terms of: including all parameters, using vague language, and parameter bloating}.

\cf{Frameworks combined} (1) \emph{Contextual integrity}: a flow is described by sender, recipient, transmission principle, attribute and subject \pc{doan2025}{6}{Using the CI framework, information flows are broken down into sender, recipient, transmission principle, attribute, and subject}; ``without specifying all five parameters, the context is too ambiguous to assess the implications'' \pc{doan2025}{6}{Without specifying all five parameters, the context is too ambiguous to assess the implications}. Operationalised following Shvartzshnaider et al.\ (2019). (2) \emph{User-relevant attributes} from Johansson et al.'s discrete-choice study on research data sharing (data collector, data user, reason for data use, information and consent, review of data sharing). (3) \emph{GDPR transparency}: Art.\ 12(1) concise/transparent/intelligible/clear language; Art.\ 13(1)a controller identity, 13(1)e recipients, 13(1)c purpose, 13(2)e whether data are a contract requirement, 13(3) secondary purposes. The mapping (Table 2) links sender $\to$ controller, recipient $\to$ recipients, transmission principle $\to$ purpose; the authors state they evaluate presence of requirements, ``we are not assessing compliance'' \pc{doan2025}{8}{While the presence or absence of relevant GDPR requirements could be systematically evaluated, we are not assessing compliance}.

\cf{Sample (stage A and D)} Six widely used DTC genetic-testing companies offering ancestry, trait and/or wellness reports, excluding medical genetic disease, paternity, microbiome and full-genome tests, and excluding business-to-business or medically supervised services; with EU/UK websites; global leaders plus EU-specific providers (23andMe, AncestryDNA, MyHeritage, FamilyTreeDNA, LivingDNA, TellMeGen) ``to try to gain a representative sample'' \pc{doan2025}{5}{We chose six of the most widely used DTC genetic testing companies with similar services of ancestry, trait, and/or wellness reports}. Policies captured from the EU or UK English websites between December 2022 and January 2023 \pc{doan2025}{5}{We accessed the EU or UK English websites’ privacy policies and research consent policies from each company from December 2022 to January 2023}; TellMeGen's differently named pages (``legal terms'', ``legal consent'') were used instead; external linked pages were excluded to limit scope \pc{doan2025}{5}{Information on external pages was not included as they were linked outside the policies}. Corpus: 10 privacy policies, 9 research consents, 1 group-project consent, 1 legal-terms document, 1 legal-consent document \pc{doan2025}{5}{our corpus included 10 privacy policies, 9 research consent}.

\cf{Procedure and coding (stage C and I)} Identify information flows about genetic data (directly or indirectly), co-code their CI parameters, then run three analyses \pc{doan2025}{6}{Even one missing parameter can be confusing to an individual because one premise of CI is that the flow must be complete to understand the full context}: \emph{missing information} (any of the five parameters absent); \emph{parameter bloating} (two or more discrete entries in a parameter, e.g.\ several purposes such as marketing, research and service provision in one transmission principle) \pc{doan2025}{7}{Parameter bloating identifies any flows with two or more discrete entries per parameter}; \emph{vague terms} by matching flow text to Bhatia et al.'s list in four lexical categories: numeric quantifiers (most, some), modals (may, likely, possibly), generalisations (mostly, normally, primarily) and conditionals (sometimes, as necessary, depending) \pc{doan2025}{8}{numeric quantifiers such as most or some, modal terms such as may, likely, and possibly}. Tools: MaxQDA and R. For RQ3, keyword search (review, ethical, opt-out, opt-in, purpose, goal) located relevant text, and coding started from Johansson's attributes with new bottom-up subcategories under ``reason for data use'' (performing a contracted service, unspecified research, secondary DNA service). Risks and benefits were found with keywords (risk, harm, benefit, compensation) and coded bottom-up, with terms taken from the corpus to ``minimize subjectivity'' \pc{doan2025}{9}{To minimize subjectivity, categories of risks and benefits were coded with terms found in the corpus}. \emph{Reliability}: coding was done by Author 1; ``in confusing cases Author 1 consulted Author 2 for higher reliability'' \pc{doan2025}{6}{While the validity of information and parameter bloating may be subjective, in confusing cases Author 1 consulted Author 2 for higher reliability}; uncertainties for user relevance were discussed, ``though not co-coded'' \pc{doan2025}{9}{Author 1 tried to be objective in adherence to Johansson et al.’s work and uncertainties were discussed with Author 2, though not co-coded}. No agreement statistic is reported. The codebook is linked in an appendix.

\cf{Results RQ1/RQ2} 62 information flows were identified \pc{doan2025}{10}{we identified 62 information flows across the companies’ publicly available privacy or consent policies directly or indirectly referencing the transfer of genetic data}; only about 30\,\% concern genetic data specifically, about 35\,\% genetic plus other personal data, about 30\,\% personal data in general. \emph{Missing information}: 10 flows (16\,\%): transmission principle 3, recipient 4, attribute 1, and one with recipient, attribute and subject missing \pc{doan2025}{10}{10 (16\%) of flows had missing information}. \emph{Bloating}: one flow had one sender, one subject, 8 attributes, 10 recipients and one broad transmission principle; across companies up to 9 attributes, 8 recipients and 14 transmission principles \pc{doan2025}{10}{Parameter bloating was found in 37\% of information flows}. \emph{Vagueness}: 81\,\% of flows had at least one vague term; modals appeared in 69\,\%, numeric quantifiers 33\,\%, conditionals 16\,\%, generalisations only for MyHeritage (8\,\%) \pc{doan2025}{12}{modal words (e.g., may, can, possibly) were used in 69\% of all surveyed companies’ information flows}. Mapping: 17\,\% of flows lack information items mandated by Art.\ 12(1), 5 flows omit the recipient (Art.\ 13(1)e), 5 lack a clear transmission principle (Art.\ 13(1)c, 13(3)). \textbf{Internal inconsistencies}: the text reports bloating in 37\,\% of flows, but Table 4 and the conclusion give 20 flows (32\,\%); the introduction mentions 59 flows while the rest says 62; the risk section reports 37 stated risks while Table 5 sums to 31. Check numbers against the tables before quoting them.

\cf{Results RQ3} Data collector: companies resemble a technological company, but it is inferred. Data user: instead of naming recipients, categories such as technological or pharmaceutical company, academic project, healthcare professionals, law enforcement, other users via genealogy services. Reason for data use: 34\,\% of flows had more than one reason; ``providing a contracted service'' 35.8\,\%, ``developing a new product or service'' 15.6\,\%, ``unspecified research'' 12.6\,\% \pc{doan2025}{13}{was the most common reason for data use at 35.8\%}. Information and consent: purposes embedded in bloated transmission principles; storage opt-in/out varied (TellMeGen destroys samples after two months, LivingDNA after ten years, FamilyTreeDNA keeps them). Review of data sharing: only 23andMe and AncestryDNA mention any ethical review, and AncestryDNA never describes its type \pc{doan2025}{14}{Only 23andMe and AncestryDNA mention any review when the data is shared in a new context}.

\cf{Results RQ4 and discussion} Stated risks: seven categories (identification incl.\ re-identification, family exposure and phenotype; undesired third parties incl.\ law enforcement and insurers; inaccurate results; secondary findings; data breaches; discomfort; general catch-all) \pc{doan2025}{14}{a total of 37 risks were identified, spanning 7 categories}; benefits fewer, mostly ``indirect data altruism'' \pc{doan2025}{15}{The number benefits is lower than risks: 5 benefits over two categories}. Privacy policies frame risk as company liability and push responsibility onto the user, while consent documents add mitigation. No policy mentions risks to genetic relatives or suggests deciding as a family \pc{doan2025}{19}{No privacy or consent policies mention the risks to genetic family members}. Interpretation: vague language may reflect completeness plus a desire to avoid frequent updates \pc{doan2025}{17}{However, companies may have more complete information but present it using vague language to cover multiple cases}; WP260 guidance recommends changing ``may'' to ``will'' where possible; algorithmic tools can pre-screen but still need expert analysis to assess compliance, usability and transparency \pc{doan2025}{17}{this would still require expert analysis to assess compliance, usability, and transparency}.

\cf{Limits (author and card)} Six companies only, possibly unrepresentative \pc{doan2025}{19}{We only assess 6 companies, which might not be very representative}; expert reading of public text only, without the companies' internal information or any contact with them; ``some results were coded by a single author and should be extended and replicated to increase reliability'', and user studies are needed to validate the findings \pc{doan2025}{20}{Some results were coded by a single author and should be extended and replicated to increase reliability}; preliminary test of the method; no inter-coder statistics; numbers inconsistent between sections; flows are counted as units with no stated unitising rule (when does a sentence contain one flow or several?); the vague-term match is a dictionary match, which yields presence, not whether vagueness harms a particular reader; no legal compliance judgement.

\cf{Use for this thesis} (1) A ready operationalisation for the policy branch of stage C: CI parameters, three flow-quality checks (missing, bloated, vague) and a published vague-term list (Bhatia et al.), which can be coded per policy statement. (2) A worked example of mapping codes to GDPR articles \emph{without} claiming legal compliance, which is the correct scoping for a bachelor thesis. (3) The weakness to avoid: one coder and an undefined flow unit. The thesis should (a) define the unit of analysis (statement / flow) and the splitting rule before coding, (b) have a second coder code a sample and report $\kappa$ or $\alpha$ per code, (c) keep dated captures with URL and version. (4) The risk/benefit categories (identification, undesired third parties, secondary findings, breach, discomfort) are a source for the harm tree in stage H, with the note that they come from genetic-testing consent documents and need adaptation to mental-health apps. (5) Transfer caveat: the study uses privacy policies of EU/UK companies, not app-store pages or in-app screens.
\end{papercard}
